\chapter{Flip-Flops}
\label{ch:ff}

So far we have built circuits \emph{without memory}: adders, comparators,
decoders. Their output is completely determined by what is on the inputs
\emph{right now}. In the security panel this already gives an alarm on the
condition ``door open and guard on'', and with a noise threshold also on
``noise above normal''. But as soon as the intruder closes the door, the
condition disappears and the siren falls silent. The question of this chapter
remains: \emph{how can a circuit remember that an alarm happened?} The answer
is an element that stores one bit---the flip-flop.

\begin{buildbox}
In Logra, open the \term{Learn} tab, the ``Flip-flops and registers'' chapter
of the ``Circuits and the board'' book, and go through the lesson ``D flip-flop:
remember a bit''. Apply a few values to the input and look at the
\term{Waveforms} tab: the output changes only at the moments of the clock edge
and keeps the past between edges. Without reading further, say in one sentence
how this differs from a circuit made of gates. The theory in this chapter will
explain how such a ``brick'' is built inside.
\end{buildbox}

\section{Why a Circuit Needs Memory}

Imagine a ``on/off'' button on a remote control: you press it and the light
comes on, you release it and it keeps burning, you press again and it goes
out. At the moment the button is released, the circuit inputs are the same, yet
the output may be either 0 or 1---depending on \emph{what happened earlier}. A
combinational circuit cannot do this. We need \term{memory}.

The simplest memory element stores one bit and is called a \term{flip-flop}
(its close relative, the latch, is described below).

\begin{diybox}[See Memory]
\textbf{Task.} \stars{1} In Logra, go through the lesson ``Remember one bit''
(section ``Time and memory'' of the Learn course). \stars{2} Open the example
\code{edge-toggle} and compare its behavior with the example \code{button}: what
happens to the output on a repeated press?\\
\textbf{In Logra:} the Learn, Code and Waveforms tabs.\\
\textbf{Expected result:} in \code{button} the output repeats the input,
in \code{edge-toggle} it also depends on the past.\\
\textbf{How to check:} in Waveforms, find two moments where the inputs are the
same but the outputs differ.
\end{diybox}

\section{The Bistable Cell}

Connect two inverters in a ring. If the output of the first is 1, then the
output of the second is 0, and this 0 keeps the output of the first at 1. The
state \emph{sustains itself}. The opposite stable state is also possible. The
circuit has two stable states---it is \term{bistable} and can store one bit.

\begin{figure}[H]
\centering
\begin{tikzpicture}
  \node[gNOT] (n1) at (0,0.9) {};
  \node[gNOT, rotate=180] (n2) at (0,-0.9) {};
  \draw[wire] (n1.output) -- ++(0.6,0) |- (n2.input); % fig-ok: feedback loop, entry segment 6 mm, rotated inverter
  \draw[wire] (n2.output) -- ++(-0.6,0) |- (n1.input);
  \node[font=\sffamily\small, text=FpgaRed] at (1.4,1.25) {1};
  \node[font=\sffamily\small, text=FpgaBlue] at (-1.75,-0.55) {0};
  \fill ($(n1.output)+(0.6,-0.9)$) circle (1.3pt);
  \draw[wire] ($(n1.output)+(0.6,-0.9)$) -- ++(0.6,0) node[right] {$Q$};
  \node[font=\sffamily\small, align=left, anchor=west] at (3.4,0.3) {The state sustains itself,};
  \node[font=\sffamily\small, align=left, anchor=west] at (3.4,-0.3) {but there is no way to change it yet.};
\end{tikzpicture}
\caption{Two inverters in a ring---a bistable cell}
\end{figure}

\begin{diybox}[A Ring of Two Inverters]
\textbf{Task.} \stars{2} In the Designer, connect two NOT blocks in a ring and
send the output of one of them to an LED. \stars{1} If you do not want to build
the loop from scratch, read the step ``One bit of memory'' in the lesson ``D
flip-flop: remember a bit'' (chapter ``Flip-flops and registers'').\\
\textbf{Expected result:} the ring holds either of the two states, but there is
no way to write a new value into it. The simulator may show an undefined value
\code{X}: the cell has no input that sets the initial state.\\
\textbf{How to check:} Waveforms (the output does not change by itself) and RTL
(two inversions in a circle).
\end{diybox}

\section{The RS Flip-Flop}

To write a new value into the cell, replace the inverters with NOR gates and
use their second inputs. The result is the \term{RS flip-flop}
(an asynchronous RS latch):
\begin{itemize}
  \item input $S$ (Set) writes 1;
  \item input $R$ (Reset) writes 0;
  \item if $S = R = 0$, the flip-flop \textbf{holds} its previous value.
\end{itemize}

\begin{figure}[H]
\centering
\begin{minipage}{0.48\textwidth}
\centering
\begin{tikzpicture}
  \node[gNOR] (g1) at (0,1.3) {};
  \node[gNOR] (g2) at (0,-1.3) {};
  \draw[wire] (g1.input 1) -- ++(-1.2,0) node[left] {$R$};
  \draw[wire] (g2.input 2) -- ++(-1.2,0) node[left] {$S$};
  \draw[wire] (g1.output) -- ++(1.4,0) node[right] {$Q$};
  \draw[wire] (g2.output) -- ++(1.4,0) node[right] {$\overline{Q}$};
  \draw[wire] ($(g1.output)+(0.5,0)$) |- (0,0.4) -| ($(g2.input 1)+(-0.45,0)$) -- (g2.input 1);
  \draw[wire] ($(g2.output)+(0.9,0)$) |- (0,-0.4) -| ($(g1.input 2)+(-0.85,0)$) -- (g1.input 2);
  \fill ($(g1.output)+(0.5,0)$) circle (1.3pt); \fill ($(g2.output)+(0.9,0)$) circle (1.3pt);
\end{tikzpicture}
\end{minipage}\hfill
\begin{minipage}{0.48\textwidth}
\centering
\begin{tabular}{cc|c|l}
\toprule
$S$ & $R$ & $Q$ & \textbf{Mode} \\
\midrule
0 & 0 & $Q$ & hold \\
1 & 0 & 1 & set \\
0 & 1 & 0 & reset \\
1 & 1 & \textcolor{FpgaRed}{?} & \textcolor{FpgaRed}{forbidden} \\
\bottomrule
\end{tabular}
\end{minipage}
\caption{An RS flip-flop made of two NOR gates and its transition table}
\label{fig:rs}
\end{figure}

Let us see why it works. Remember: the output of a NOR gate is 0 if at least
one input is 1.
\begin{itemize}
  \item $S=1$: the lower gate outputs $\overline{Q}=0$. Both inputs of the upper
        gate are zero ($R=0$, $\overline{Q}=0$), so $Q=1$.
  \item Now $S$ returns to 0. On the lower gate $Q=1$, so $\overline{Q}$ is
        still 0, and $Q$ is still 1: the flip-flop has
        \emph{remembered} the one.
  \item $S=R=1$---both outputs become zeros (although they should be
        opposite), and when the signals are released at the same time, it is
        unknown which state the flip-flop will ``fall'' into. That is why this
        combination is forbidden.
\end{itemize}

\begin{figure}[H]
\centering
\begin{tikztimingtable}[timing/slope=0.1, xscale=3.0, yscale=1.4, timing/font=\sffamily\small, timing/rowdist=3.2ex]
  $S$ & 1L 1H 10L \\
  $R$ & 5L 1H 6L \\
  $Q$ & 1L 4H 7L \\
\extracode
  \node[lbl, text=FpgaOrange] at (1.5,1.7) {set};
  \node[lbl, text=FpgaOrange] at (3.5,1.7) {hold 1};
  \node[lbl, text=FpgaOrange] at (5.5,1.7) {reset};
  \node[lbl, text=FpgaOrange] at (8.0,1.7) {hold 0};
\end{tikztimingtable}
\caption{Operation of the RS flip-flop: a short pulse on $S$ or $R$ switches the
flip-flop, and without pulses it holds its value}
\end{figure}

\begin{infobox}[An RS Flip-Flop Made of NAND Gates]
The flip-flop can be built in exactly the same way from two NAND gates. But then
the inputs act on \emph{zero}: $\overline{S}=0$ sets, $\overline{R}=0$
resets, and the combination $\overline{S}=\overline{R}=0$ is forbidden. Such a
flip-flop is a classic way to suppress the contact bounce of a mechanical
switch: at the very first touch of the contact it switches and does not notice
the further ``jumps''.
\end{infobox}

\begin{diybox}[The RS Flip-Flop]
\textbf{Task.} \stars{1} The lesson ``SR flip-flop: set and reset''; it also
has a step ``The forbidden combination''. \stars{2} Apply $S=R=1$, then release
both signals and describe what happened. \stars{3} Build a ``light-switch
button'' from two buttons: one turns the light on, the other turns it off.\\
\textbf{In Logra:} Learn, Designer, Waveforms.\\
\textbf{Expected result:} a short press of $S$ turns the light on, and it stays
on after release; pressing both at once gives an ambiguous result.\\
\textbf{How to check:} Waveforms: run the sequence $S$, pause,
$R$, pause and compare with Fig.~\ref{fig:rs}.
\end{diybox}

\subsection{The RS Flip-Flop in Verilog}
\label{sec:ff-sr-v}

The same circuit can be described in Verilog without drawing anything with the
mouse: each NOR gate becomes one \code{assign} line. The output of each gate
depends on the output of the other one---this is the feedback loop.

\vfile{\codepath/b7_flipflops/sr_latch.v}{code/b7\_flipflops/sr\_latch.v}

\textbf{How to read it.} The first \code{assign} is the upper NOR gate:
$Q=\overline{R+\overline{Q}}$. The second is the lower one:
$\overline{Q}=\overline{S+Q}$. The keyword \code{assign} describes \emph{wires
and gates}, not a sequence of actions, so the order of the lines does not
matter, and the fact that \code{q} and \code{qn} refer to each other is exactly
what gives memory. The testbench \code{tb\_sr\_latch.v} (next to the module in
the folder \code{code-en/b7\_flipflops}) makes four steps: $S=1$, then $S=0$
(holding the one), $R=1$, then $R=0$ (holding the zero). A value has to be
written first ($S=1$ first): until the flip-flop has stored something, its
output is unknown (\code{X}), and that is normal.

\textbf{How to run it in Logra.} Open the \term{Code} tab, paste the module and
press \term{Synthesize} (Ctrl+Enter). Logra warns that a logic loop was found
in the module: for memory built from gates this is expected, whereas in an
ordinary combinational circuit such a warning would mean a mistake. On the
\term{RTL} tab you will see two NOR gates closed in a ring; in \term{Waveforms}
apply the sequence $S$, pause, $R$, pause and compare with
Fig.~\ref{fig:rs}.

\begin{diybox}[RS in Verilog: Complete the Module]
\textbf{Task.} \stars{2} Describe an RS flip-flop made of NAND gates (the inputs
are active-low: \code{sn} sets, \code{rn} resets). \stars{3} Find out in
Waveforms what happens at \code{sn = rn = 0}, and then when both go to
\code{1} at once (this is forbidden for NOR, but not in this circuit!). How do
the forbidden combinations of the two versions differ?\\
\textbf{Expected result:} \code{sn = 0} gives \code{q = 1}, \code{rn = 0} gives
\code{q = 0}, and at \code{sn = rn = 1} the circuit holds.
\end{diybox}

\textbf{Sample answer.}
\begin{vcode}
module sr_latch_n (
    input  sn,       // set, active level 0
    input  rn,       // reset, active level 0
    output q,
    output qn
);
    assign q  = ~(sn & qn);   // NAND
    assign qn = ~(rn & q);    // NAND
endmodule
\end{vcode}
At \code{sn = rn = 0} both outputs become ones (the forbidden combination of
this version); after both signals are released at the same time the state is
unknown.

\section{The D Latch}

The RS flip-flop has two inconveniences: there is a forbidden input combination,
and it reacts to the inputs at any moment. Put two AND gates in front of it,
a data input $D$ and an enable input $C$ (clock). The result is the
\term{D latch}:
\begin{itemize}
  \item when $C=1$ the latch is \textbf{transparent}: the output follows the
        input, $Q = D$;
  \item when $C=0$ the latch \textbf{holds} the value that $D$ had at the moment
        $C$ went from 1 to 0.
\end{itemize}

\begin{figure}[H]
\centering
\begin{tikzpicture}
  \node[gAND] (a1) at (3,1.2) {};
  \node[gAND] (a2) at (3,-1.2) {};
  \node[draw=FpgaGreen, fill=green!8, thick, minimum width=18mm, minimum height=36mm, font=\sffamily\small, align=center] (rs) at (6.6,0) {RS\\flip-flop};
  \node[font=\sffamily\footnotesize] at ($(rs.west)+(0.25,1.2)$) {R};
  \node[font=\sffamily\footnotesize] at ($(rs.west)+(0.25,-1.2)$) {S};
  \node[gNOT, anchor=output] (n) at ($(a1.input 1)+(-0.9,0)$) {};
  \coordinate (Dp) at ($(a2.input 2)+(-3.0,0)$);
  \coordinate (Dt) at ($(Dp)+(0.5,0)$);
  \draw[wire] (Dp) node[left] {$D$} -- (a2.input 2);
  \fill (Dt) circle (1.3pt);
  \draw[wire] (Dt) -- (Dt |- n.input) -- (n.input);
  \draw[wire] (n.output) -- (a1.input 1);
  \coordinate (cb1) at ($(a1.input 2)+(-0.4,0)$);
  \coordinate (cb2) at ($(a2.input 1)+(-0.4,0)$);
  \draw[wire, FpgaBlue] ($(Dp |- 0,0)$) node[left, text=FpgaDark] {$C$} -- (cb1 |- 0,0);
  \draw[wire, FpgaBlue] (cb1 |- 0,0) -- (cb1) -- (a1.input 2);
  \draw[wire, FpgaBlue] (cb1 |- 0,0) -- (cb2) -- (a2.input 1);
  \fill[FpgaBlue] (cb1 |- 0,0) circle (1.3pt);
  \draw[wire] (a1.output) -- ($(rs.west)+(0,1.2)$);
  \draw[wire] (a2.output) -- ($(rs.west)+(0,-1.2)$);
  \draw[wire] ($(rs.east)+(0,1.2)$) -- ++(0.8,0) node[right] {$Q$};
  \draw[wire] ($(rs.east)+(0,-1.2)$) -- ++(0.8,0) node[right] {$\overline{Q}$};
\end{tikzpicture}
\caption{The D latch: the forbidden combination $S=R=1$ is now impossible,
because $S$ and $R$ receive $D$ and $\overline{D}$}
\end{figure}

\begin{figure}[H]
\centering
\begin{tikztimingtable}[timing/slope=0.08, xscale=1.6, yscale=1.4, timing/font=\sffamily\small, timing/rowdist=3.2ex]
  $C$ & 2L 4H 4L 4H 2L \\
  $D$ & 1L 2H 1L 4H 3L 1H 4L \\
  $Q$ & 2L 1H 1L 6H 1L 1H 4L \\
\extracode
  \begin{pgfonlayer}{background}
    \fill[green!10] (2,0.6) rectangle (6,-5.2);
    \fill[green!10] (10,0.6) rectangle (14,-5.2);
  \end{pgfonlayer}
  \node[lbl, text=FpgaGreen] at (4,1.7) {transparent};
  \node[lbl] at (8,1.7) {holds};
  \node[lbl, text=FpgaGreen] at (12,1.7) {transparent};
\end{tikztimingtable}
\caption{At $C=1$ the D latch passes all changes of $D$ to the output; at $C=0$
it ``freezes'' the output}
\label{fig:dlatch_t}
\end{figure}

\begin{diybox}[The D Latch]
\textbf{Task.} \stars{2} In the Designer, build the latch of Fig.~\ref{fig:dlatch_t}
from gates: NOT, two ANDs and an RS flip-flop made of two NORs. \stars{3} Apply
a changing $D$ and the enable $C$ as a long pulse.\\
\textbf{In Logra:} Designer, Waveforms.\\
\textbf{Expected result:} while $C=1$ the output follows $D$ (including all its
changes), and at $C=0$ the value is frozen.\\
\textbf{How to check:} Waveforms; compare with Fig.~\ref{fig:dlatch_t}.
Save the circuit: you will need it in the next section.
\end{diybox}

\subsection{The Latch in Verilog: \code{always @*}}
\label{sec:ff-latch-v}

In Verilog a latch is described by an \code{always @*} process: a combinational
block that ``wakes up'' whenever any of the inputs it reads changes. If not all
branches of the block assign a value to the output, the synthesizer concludes
that under the remaining conditions the output must be \emph{remembered}, and
inserts a latch.

\vfile{\codepath/b7_flipflops/d_latch.v}{code/b7\_flipflops/d\_latch.v}

\textbf{How to read it.} ``Whenever \code{c} or \code{d} changes: if \code{c} is
1, then \code{q} equals \code{d}.'' There is no \code{else} branch: at
\code{c = 0} nothing is assigned to the output, which means it holds the old
value. That is how memory appears. Here the latch is intended, and its
assignment is \emph{blocking} (\code{=}), as is customary in combinational
blocks.

\textbf{How to run it in Logra.} \term{Code} tab, Ctrl+Enter. Logra warns about
an inferred latch (in the Yosys log: \emph{Latch inferred}). In a teaching
module this confirms that you got what you wanted; in a working project such a
message is a reason to check whether you forgot an \code{else} branch. On the
\term{RTL} tab find the latch block; in \term{Waveforms} apply $D$ as fast
pulses during a long $C=1$ and watch $Q$ follow $D$.

\begin{diybox}[The Latch in Verilog: Change the Module]
\textbf{Task.} \stars{1} Change the latch so that it is transparent at $c=0$
rather than $c=1$. \stars{2} Remove the condition from the module and write
\code{q = d} unconditionally: what has the circuit turned into, and what does
Logra say now?\\
\textbf{Expected result:} the latch warning disappears, and the circuit is just a
wire from \code{d} to \code{q}.
\end{diybox}

\textbf{Sample answer.}
\begin{vcode}
module d_latch_low (input c, d, output reg q);
    always @* begin
        if (!c) q = d;       // transparent at c = 0
    end
endmodule
\end{vcode}
Without the condition we get \code{always @* q = d;}: this is a combinational
circuit without memory, and the latch warning goes away.

\section{The D Flip-Flop}

The transparency of a latch is both an advantage and a drawback. If the output
of a latch goes back to its own input through combinational logic (as in a
counter), then while $C=1$ the signal will run around the loop, and the result
will be unpredictable. We need an element that stores the value at \emph{one
moment in time}---at the \emph{edge} of the clock signal. This is the
\term{D flip-flop}, the main memory element of all modern digital
technology.

It can be built from two D latches connected in series with opposite enable
signals. Such a circuit is called \term{master--slave}:

\begin{figure}[H]
\centering
\begin{tikzpicture}[lt/.style={draw=FpgaGreen, fill=green!8, thick, minimum width=24mm, minimum height=18mm, font=\sffamily\small, align=center}]
  \node[lt] (m) at (0,0) {\quad master\\\quad latch};
  \node[lt] (s) at (4.2,0) {\quad slave\\\quad latch};
  \node[font=\sffamily\footnotesize] at ($(m.west)+(0.2,0.45)$) {D};
  \node[font=\sffamily\footnotesize] at ($(m.west)+(0.2,-0.45)$) {C};
  \node[font=\sffamily\footnotesize] at ($(s.west)+(0.2,0.45)$) {D};
  \node[font=\sffamily\footnotesize] at ($(s.west)+(0.2,-0.45)$) {C};
  \draw[wire] ($(m.west)+(-1.2,0.45)$) node[left] {$D$} -- ($(m.west)+(0,0.45)$);
  \draw[wire] ($(m.east)+(0,0.45)$) -- node[above, lbl] {$Q_m$} ($(s.west)+(0,0.45)$);
  \draw[wire] ($(s.east)+(0,0.45)$) -- ++(1,0) node[right] {$Q$};
  \node[gNOT, anchor=output] (n) at ($(m.west)+(-0.5,-0.45)$) {};
  \coordinate (ck) at ($(n.input)+(-0.5,0)$);
  \draw[wire, FpgaBlue] ($(ck)+(-0.4,0)$) node[left, text=FpgaDark] {CLK} -- (n.input);
  \draw[wire, FpgaBlue] (n.output) -- ($(m.west)+(0,-0.45)$);
  \draw[wire, FpgaBlue] (ck) |- (2.4,-1.9) |- ($(s.west)+(0,-0.45)$);
  \fill[FpgaBlue] (ck) circle (1.3pt);
  \node[font=\sffamily\small, align=left, anchor=north west] at (-2.4,-2.2) {CLK $=0$: master transparent, slave holds;\\
  CLK $=1$: master holds, slave transparent;\\ $\Rightarrow$ $Q$ changes only when CLK goes from 0 to 1.};
\end{tikzpicture}
\caption{A D flip-flop built as a master--slave circuit}
\label{fig:ms}
\end{figure}

In schematics a D flip-flop is drawn as a rectangle whose clock input is
marked with a triangle (the ``edge-triggered'' sign):

\begin{figure}[H]
\centering
\begin{tikzpicture}
  \draw[thick, fill=green!8] (0,0) rectangle (1.8,2.2);
  \node[font=\sffamily\small] at (0.3,1.7) {D}; \node[font=\sffamily\small] at (1.5,1.7) {Q};
  \node[font=\sffamily\small] at (1.5,0.5) {$\overline{\text{Q}}$};
  \draw[thick] (0,0.3) -- (0.25,0.5) -- (0,0.7);
  \draw[wire] (-0.7,1.7) node[left] {$D$} -- (0,1.7);
  \draw[wire] (-0.7,0.5) node[left] {CLK} -- (0,0.5);
  \draw[wire] (1.8,1.7) -- (2.5,1.7) node[right] {$Q$};
  \draw[wire] (1.8,0.5) -- (2.1,0.5); \draw[thick] (2.18,0.5) circle (0.08); \draw[wire] (2.26,0.5) -- (2.5,0.5) node[right] {$\overline{Q}$};
  \draw[wire] (0.9,2.2) -- (0.9,2.7) node[above] {$\overline{\text{SET}}$}; \draw[thick, fill=white] (0.9,2.28) circle (0.08);
  \draw[wire] (0.9,0) -- (0.9,-0.5) node[below] {$\overline{\text{RST}}$}; \draw[thick, fill=white] (0.9,-0.08) circle (0.08);
  \node[font=\sffamily\small, align=left, anchor=west] at (4.2,1.6) {triangle---triggers on the edge};
  \node[font=\sffamily\small, align=left, anchor=west] at (4.2,1.0) {circle---active level is 0};
  \node[font=\sffamily\small, align=left, anchor=west] at (4.2,0.4) {SET, RST---asynchronous set and reset:};
  \node[font=\sffamily\small, align=left, anchor=west] at (4.2,-0.1) {they act immediately, without waiting for an edge};
\end{tikzpicture}
\caption{The schematic symbol of a D flip-flop}
\end{figure}

\begin{figure}[H]
\centering
\begin{tikztimingtable}[timing/slope=0.08, xscale=1.6, yscale=1.4, timing/font=\sffamily\small, timing/rowdist=3.2ex]
  CLK                  & 2L 4H 4L 4H 2L \\
  $D$                  & 1L 2H 1L 4H 3L 1H 4L \\
  $Q$ of the latch     & 2L 1H 1L 6H 1L 1H 4L \\
  $Q$ of the flip-flop & 2L 8H 6L \\
\extracode
  \begin{pgfonlayer}{background}
    \draw[FpgaOrange, dashed] (2,1.3) -- (2,-7.2);
    \draw[FpgaOrange, dashed] (10,1.3) -- (10,-7.2);
  \end{pgfonlayer}
  \node[lbl, text=FpgaOrange] at (2,1.6) {edge}; \node[lbl, text=FpgaOrange] at (10,1.6) {edge};
\end{tikztimingtable}
\caption{Comparison: the latch (Fig.~\ref{fig:dlatch_t}) passes all changes of
$D$ while CLK $=1$, whereas the flip-flop takes the value of $D$ \emph{only at the edge}}
\label{fig:latch_vs_ff}
\end{figure}

\begin{diybox}[Latch Versus Flip-Flop]
\textbf{Task.} \stars{1} The lesson ``D flip-flop: remember a bit''. \stars{2} Apply
the same signal $D$ to the latch from the previous section and to the ``D
flip-flop'' block, and make the clock a long pulse. Find the moment where their
outputs diverge.
\stars{3} Apply $D=\overline{Q}$ first to the flip-flop, then to the latch.\\
\textbf{In Logra:} Designer, Waveforms, RTL.\\
\textbf{Expected result:} the flip-flop divides the frequency by 2, and the latch
at $C=1$ goes into fast oscillations (in the simulator this appears as \code{X} or
alternation).\\
\textbf{How to check:} Waveforms; compare with Fig.~\ref{fig:latch_vs_ff}.
\end{diybox}

\subsection{The Flip-Flop in Verilog: \code{always @(posedge clk)}}
\label{sec:ff-dff-v}

A D flip-flop is described by a process that wakes up \emph{only at the rising
edge} of the clock signal. Compare it with the latch: only one line changes, but
the behavior is fundamentally different.

\vfile{\codepath/b7_flipflops/d_ff.v}{code/b7\_flipflops/d\_ff.v}

\textbf{How to read it.} ``At every rising edge of \code{clk}, store in \code{q}
what is on \code{d} right now.'' The assignment here is \emph{non-blocking}
(\code{<=}): all flip-flops of the circuit take the old values at the same time
and update them at the same time. A rule that will save you from half of all
mistakes: in \code{always @(posedge ...)} use only \code{<=}, in \code{always @*}
use only \code{=}. We will see why in Chapter~\ref{ch:regs}, where a bit ``flies''
through the whole shift register in one clock cycle if the rule is broken.

The testbench \code{tb\_latch\_ff.v} applies the same signal $D$ to the latch
and to the flip-flop:

\vfile{\codepath/b7_flipflops/tb_latch_ff.v}{code/b7\_flipflops/tb\_latch\_ff.v}

\textbf{How to read the testbench.} The task \code{chk} waits 1~ns after a change
and compares both outputs with the expected ones; the operator \code{!==} counts
an unknown \code{X} as a difference, not as a match. At the end the testbench
prints \code{OK} or the number of errors.

\textbf{How to run it in Logra.} \term{Code} tab: paste \code{d\_ff} (and
\code{d\_latch} separately) and press Ctrl+Enter. A ready-made pair ``latch and
flip-flop in one module'' with common inputs is the module \code{dff\_vs\_latch}
of Chapter~\ref{ch:bridge}. In \term{Waveforms} apply a long pulse to
\code{clk} and change \code{d} inside it. The latch output changes together with
\code{d}, the flip-flop output only at the edge
(Fig.~\ref{fig:latch_vs_ff}).

\begin{diybox}[The Flip-Flop in Verilog: Change the Module]
\textbf{Task.} \stars{1} Make a flip-flop that acts on the \emph{falling} edge
(from 1 to 0). \stars{2} Add a second output \code{qn} that is always
$\overline{Q}$. \stars{3} Build a master--slave flip-flop from two such modules
(Fig.~\ref{fig:ms}) and compare it with \code{d\_ff}.\\
\textbf{Expected result:} the output changes at the falling edge of \code{clk}.
\end{diybox}

\textbf{Sample answer.}
\begin{vcode}
module d_ff_neg (
    input      clk,
    input      d,
    output reg q,
    output     qn
);
    always @(negedge clk) q <= d;   // negedge: falling edge
    assign qn = ~q;
endmodule
\end{vcode}

\section{Flip-Flop Timing Parameters}

For a flip-flop to store a value reliably, the input $D$ must stay unchanged for
some time before the edge and after it:
\begin{itemize}
  \item $t_\text{su}$ (setup, the \term{setup time})---how long $D$ must
        be stable \emph{before} the edge;
  \item $t_\text{h}$ (hold, the \term{hold time})---how long \emph{after} the edge;
  \item $t_\text{co}$ (clock-to-output)---how long after the edge $Q$ changes.
\end{itemize}

\begin{figure}[H]
\centering
\begin{tikzpicture}[x=1cm, y=1cm]
  \draw[very thick, FpgaBlue] (0,2.2) -- (4,2.2) -- (4.1,3) -- (8,3);
  \node[font=\sffamily\small, anchor=east] at (-0.1,2.6) {CLK};
  \draw[very thick] (0,1.1) -- (1.5,1.1) -- (1.7,1.5) -- (6,1.5) -- (6.2,1.1) -- (8,1.1);
  \draw[very thick] (0,1.5) -- (1.5,1.5) -- (1.7,1.1) -- (6,1.1) -- (6.2,1.5) -- (8,1.5);
  \node[font=\sffamily\small, anchor=east] at (-0.1,1.3) {$D$};
  \node[font=\sffamily\footnotesize] at (3.9,1.3) {stable};
  \draw[very thick, FpgaGreen] (0,0) -- (4.8,0) -- (4.9,0.4) -- (8,0.4);
  \node[font=\sffamily\small, anchor=east] at (-0.1,0.2) {$Q$};
  \fill[FpgaRed, opacity=0.15] (2.6,0.9) rectangle (4.05,1.7);
  \fill[FpgaOrange, opacity=0.2] (4.05,0.9) rectangle (4.8,1.7);
  \draw[dashed, FpgaGray] (4.05,3.3) -- (4.05,-0.5);
  \draw[<->] (2.6,-0.35) -- (4.05,-0.35) node[midway, below, font=\sffamily\small] {$t_\text{su}$};
  \draw[<->] (4.05,-0.35) -- (4.8,-0.35) node[midway, below, font=\sffamily\small] {$t_\text{h}$};
  \draw[<->] (4.05,-1.1) -- (4.85,-1.1) node[midway, below, font=\sffamily\small] {$t_\text{co}$};
\end{tikzpicture}
\caption{Setup time, hold time and the output delay of a flip-flop}
\label{fig:setup_hold}
\end{figure}

\begin{warnbox}[Metastability]
If $D$ changes inside the window $t_\text{su}$--$t_\text{h}$, the flip-flop may
``hang'' for some time in an intermediate state---neither 0 nor 1---and
only then fall randomly into one of them. This is \term{metastability}. It
inevitably arises when a signal unrelated to the clock (for example, from a
button) arrives at a flip-flop. The protection is a \term{synchronizer} made of
two flip-flops in a row: by the time of the next edge the first flip-flop has
almost certainly settled.
\end{warnbox}

\begin{diybox}[An Asynchronous Button]
\textbf{Task.} \stars{1} The lesson ``Synchronizer: tie the button to the clock''.
\stars{2} In Waveforms, move the button press to different moments of the clock
and explain on which clock cycle the result becomes visible. \stars{3} Add the
``Edge detector: one pulse per press'' (a lesson of the same name) after the
synchronizer.\\
\textbf{In Logra:} Learn, Designer, Waveforms.\\
\textbf{Expected result:} the synchronizer output changes only on an edge,
lagging by two clock cycles; the detector gives exactly one pulse per press,
however many cycles the button is held.\\
\textbf{How to check:} Waveforms: holding the button for 10 cycles gives
one pulse.
\end{diybox}

\subsection{Synchronizer and Edge Detector in Verilog}
\label{sec:ff-sync-v}

A synchronizer is exactly two flip-flops in a row, and in Verilog it takes three
lines:

\vfile{\codepath/b7_flipflops/sync2.v}{code/b7\_flipflops/sync2.v}

\textbf{How to read it.} Both assignments are in one block and both are
non-blocking: at one edge \code{ff1} gets the input value, while \code{ff2} gets
the \emph{old} value of \code{ff1}. So the value passes through both flip-flops
in two clock cycles, not one. Metastability can occur only in \code{ff1}, and it
has a whole clock cycle to die out before \code{ff2}.

An edge detector remembers the previous value of the signal and compares it with
the current one:

\vfile{\codepath/b7_flipflops/edge_det.v}{code/b7\_flipflops/edge\_det.v}

\textbf{How to read it.} \code{prev} is the value of \code{in} in the previous
clock cycle. The expression \code{in \& \textasciitilde prev} is true when it is
1 now and was 0 before: one clock cycle per press. This is a \emph{combinational}
output, so it is written with \code{assign}, not in a clocked block.

The chain ``button $\to$ synchronizer $\to$ edge detector'' is collected in one
module:

\vfile{\codepath/b7_flipflops/button_pulse.v}{code/b7\_flipflops/button\_pulse.v}

The testbench \code{tb\_button\_pulse.v} presses the button \emph{away from the
clock edge}, holds it for 10 cycles, releases it and counts the pulses: there
must be exactly one. It also checks that \code{sync2} answers not in one clock
cycle but in two.

\textbf{How to run it in Logra.} \term{Code} tab: paste \code{button\_pulse},
Ctrl+Enter, \term{RTL}: three flip-flops and one AND gate with an inverted
input. In \term{Waveforms} apply a long pulse of 10 cycles to \code{btn} and
see that \code{pulse} lasts one clock cycle and lags the button by two. For the
board the button is connected to a pin in \term{Pins}; do not forget that a
mechanical button also bounces (Chapter~\ref{ch:examples}).

\begin{diybox}[Edge Detector: Change the Module]
\textbf{Task.} \stars{1} Make a detector of the \emph{falling} edge (button
release). \stars{2} Make a detector of \emph{any} change of the signal (both
edges). \stars{3} Build a chain: a button press toggles an LED (synchronizer,
edge detector, T flip-flop).\\
\textbf{Expected result:} one pulse per release; two pulses per press and
release; the LED changes only on the press.
\end{diybox}

\textbf{Sample answer.}
\begin{vcode}
module edge_fall (input clk, in, output pulse);
    reg prev = 1'b0;
    always @(posedge clk) prev <= in;
    assign pulse = ~in & prev;       // 0 now, 1 before
endmodule

module edge_any (input clk, in, output pulse);
    reg prev = 1'b0;
    always @(posedge clk) prev <= in;
    assign pulse = in ^ prev;        // changed in either direction
endmodule

module button_toggle (input clk, btn, output reg led);
    reg ff1 = 1'b0, ff2 = 1'b0, ff3 = 1'b0;
    initial led = 1'b0;
    always @(posedge clk) begin
        ff1 <= btn; ff2 <= ff1; ff3 <= ff2;
        if (ff2 & ~ff3) led <= ~led;   // the pulse toggles the T flip-flop
    end
endmodule
\end{vcode}

\section{The Enable Input}

Often we need a flip-flop to update not on every edge, but only on command. For
this a clock enable input CE is added. It is simple: a multiplexer is placed in
front of the input $D$, and at CE $=0$ it feeds the flip-flop's own output back
to its input---the flip-flop rewrites itself and does not change.

\begin{figure}[H]
\centering
\begin{tikzpicture}
  \draw[thick, fill=orange!15] (0,0.9) -- (0.9,0.5) -- (0.9,-0.5) -- (0,-0.9) -- cycle;
  \node[font=\sffamily\footnotesize] at (0.45,0) {MUX};
  \node[font=\footnotesize, anchor=west] at (0.02,0.5) {0}; \node[font=\footnotesize, anchor=west] at (0.02,-0.5) {1};
  \draw[thick, fill=green!8] (1.8,-0.8) rectangle (3.2,0.8);
  \node[font=\sffamily\small] at (2.05,0) {D}; \node[font=\sffamily\small] at (2.95,0) {Q};
  \draw[thick] (2.3,-0.8) -- (2.5,-0.6) -- (2.7,-0.8);
  \draw[wire] (0.9,0) -- (1.8,0);
  \draw[wire] (3.2,0) -- (4.4,0) node[right] {$Q$};
  \draw[wire] (3.8,0) -- (3.8,1.5) -| (-0.6,0.5) -- (0,0.5); \fill (3.8,0) circle (1.3pt);
  \draw[wire] (-1.2,-0.5) node[left] {$D$} -- (0,-0.5);
  \draw[wire, FpgaBlue] (0.45,-0.7) -- (0.45,-1.4) node[below, text=FpgaDark] {CE};
  \draw[wire, FpgaBlue] (2.5,-0.8) -- (2.5,-1.4) node[below, text=FpgaDark] {CLK};
\end{tikzpicture}
\caption{A D flip-flop with an enable input: a multiplexer plus feedback}
\label{fig:ce}
\end{figure}

\begin{diybox}[Write Enable]
\textbf{Task.} \stars{1} The lesson ``A register with write enable''. \stars{2}
The lesson ``A register with reset'': add a synchronous reset and check that
when reset and write happen together, the reset wins. \stars{3} Build the same
without the ready-made block: a multiplexer plus the ``D flip-flop'' block as in
Fig.~\ref{fig:ce}.\\
\textbf{In Logra:} Learn, Designer, Waveforms, RTL.\\
\textbf{Expected result:} at $\text{CE}=0$ the output holds its value, however
$D$ changes.\\
\textbf{How to check:} Waveforms (change $D$ at $\text{CE}=0$); RTL
will show the same multiplexer in front of the flip-flop.
\end{diybox}

\subsection{Enable and Reset in Verilog}
\label{sec:ff-en-rst-v}

In Verilog the enable input is simply an \code{if} inside a clocked block. The
synthesizer builds the multiplexer with feedback (Fig.~\ref{fig:ce}) by itself.

\vfile{\codepath/b7_flipflops/dff_en.v}{code/b7\_flipflops/dff\_en.v}

\textbf{How to read it.} ``At the edge: if \code{en} is 1, store \code{d}.''
There is no \code{else} branch, and that is \emph{correct}: in a clocked block
the absence of an assignment means ``\code{q} stays as it was'', that is, the
flip-flop rewrites itself. The declaration \code{output reg q = 1'b0} sets the
initial value after power-up.

The testbench \code{tb\_dff\_en.v} checks three clock cycles: \code{en=0}---the
value does not change; \code{en=1}---it is stored; \code{en=0} again---it is
held. In Logra: the \term{Code} tab, Ctrl+Enter, \term{RTL} (a multiplexer in
front of the flip-flop), \term{Waveforms}: keep \code{en=0} and change \code{d},
then apply a pulse to \code{en}.

Now the \term{reset}. It comes in two kinds, and the difference is visible
already in the first line of the block:

\vfile{\codepath/b7_flipflops/dff_rst_sync.v}{code/b7\_flipflops/dff\_rst\_sync.v}

\vfile{\codepath/b7_flipflops/dff_rst_async.v}{code/b7\_flipflops/dff\_rst\_async.v}

\textbf{How to read it.} The \emph{synchronous} reset has only \code{posedge clk}
in its sensitivity list: \code{rst} is checked inside the block, that is, only at
the moment of the edge, and it acts like an ordinary data input (the flip-flop
gets an extra gate at its input \code{d}). The \emph{asynchronous} one has
\code{or posedge rst} in the list: the flip-flop is cleared immediately when
\code{rst} rises, without waiting for the clock. An FPGA flip-flop has a
separate asynchronous reset input (like RST on the symbol above), so no extra
logic is added on \code{d}. The order of the branches
\code{if (rst) \ldots\ else \ldots} matters: the reset is written first, so it
is stronger than the write.

\vfile{\codepath/b7_flipflops/tb_dff_rst.v}{code/b7\_flipflops/tb\_dff\_rst.v}

The testbench applies the reset \emph{between} edges: the synchronous
flip-flop does not notice it until the next edge, the asynchronous one is
cleared at once. In Logra repeat it in \term{Waveforms}: build both modules,
keep \code{d=1}, raise \code{rst} between edges and compare the moments when the
output fell to 0.

\begin{diybox}[Reset and Enable: Complete the Module]
\textbf{Task.} \stars{1} Add to \code{dff\_en} a synchronous reset that is
stronger than the write. \stars{2} Make an asynchronous \emph{set} to 1
(\code{set}) instead of a reset. \stars{3} Make an asynchronous reset with
active level 0 (input \code{rst\_n}): the sensitivity list will need
\code{negedge}.\\
\textbf{Expected result:} at \code{rst = 1} and \code{en = 1} the output is 0.
\end{diybox}

\textbf{Sample answer.}
\begin{vcode}
module dff_en_rst (
    input      clk, rst, en, d,
    output reg q
);
    always @(posedge clk) begin
        if (rst)     q <= 1'b0;    // reset first: it is stronger
        else if (en) q <= d;
    end
endmodule

module dff_aset (
    input      clk, set, d,
    output reg q
);
    always @(posedge clk or posedge set) begin
        if (set) q <= 1'b1;        // asynchronous set to 1
        else     q <= d;
    end
endmodule

module dff_arst_n (
    input      clk, rst_n, d,
    output reg q
);
    always @(posedge clk or negedge rst_n) begin   // reset at the fall of rst_n
        if (!rst_n) q <= 1'b0;                      // active level 0
        else        q <= d;
    end
endmodule
\end{vcode}
The entry in the list and the check inside the block must agree: \code{negedge
rst\_n} in the list means \code{if (!rst\_n)} inside.

\section{The T Flip-Flop and the JK Flip-Flop}

The \term{T flip-flop} (toggle) changes its state to the opposite one at every
edge if the input is $T=1$, and holds it at $T=0$. It is easy to make from a D
flip-flop and an XOR gate: $D = T \oplus Q$.

\begin{figure}[H]
\centering
\begin{minipage}{0.5\textwidth}
\centering
\begin{tikzpicture}
  \node[gXOR] (x) at (0,0) {};
  \draw[thick, fill=green!8] (1.8,-0.8) rectangle (3.2,0.8);
  \node[font=\sffamily\small] at (2.05,0) {D}; \node[font=\sffamily\small] at (2.95,0) {Q};
  \draw[thick] (2.3,-0.8) -- (2.5,-0.6) -- (2.7,-0.8);
  \draw[wire] (x.output) -- (1.8,0);
  \draw[wire] (3.2,0) -- (4.1,0) node[right] {$Q$};
  \draw[wire] (3.65,0) -- (3.65,1.4) -| ($(x.input 1)+(-0.5,0)$) -- (x.input 1); \fill (3.65,0) circle (1.3pt);
  \draw[wire] (x.input 2) -- ++(-0.9,0) node[left] {$T$};
  \draw[wire, FpgaBlue] (2.5,-0.8) -- (2.5,-1.3) node[below, text=FpgaDark] {CLK};
\end{tikzpicture}
\end{minipage}\hfill
\begin{minipage}{0.45\textwidth}
\centering
\begin{tabular}{c|c|l}
\toprule
$T$ & $Q$ after the edge & \\
\midrule
0 & $Q$ & hold \\
1 & $\overline{Q}$ & toggle \\
\bottomrule
\end{tabular}
\end{minipage}
\caption{A T flip-flop made of a D flip-flop and an XOR gate}
\end{figure}

If you keep $T=1$ on a T flip-flop, it toggles at every edge, and the frequency
at its output is \emph{half} the clock frequency. This is the simplest
\term{divide-by-2 frequency divider}.

\begin{figure}[H]
\centering
\begin{tikztimingtable}[timing/slope=0.08, xscale=1.5, yscale=1.4, timing/font=\sffamily\small, timing/rowdist=3.2ex]
  CLK            & 8{1L 1H} \\
  $Q$ ($T=1$)    & 1L 2H 2L 2H 2L 2H 2L 2H 1L \\
\end{tikztimingtable}
\caption{A T flip-flop with $T=1$ divides the frequency by 2}
\end{figure}

\textbf{The T flip-flop in Verilog.} The condition ``if $T=1$, toggle'' is one
\code{if} and an inversion:

\vfile{\codepath/b7_flipflops/t_ff.v}{code/b7\_flipflops/t\_ff.v}

\textbf{How to read it.} At the edge: if reset, write 0; otherwise, if \code{t} is
1, write the \emph{inversion} of itself (\code{~q}). The statement
\code{q <= ~q} reads the old value of \code{q} and writes the new one, so it is
not a contradiction but exactly ``$D=T\oplus Q$''. If you leave \code{t = 1},
you get a divide-by-2 frequency divider. The testbench \code{tb\_t\_ff.v} checks
that at $t=0$ the value is held and at $t=1$ it changes at every edge. In Logra:
\term{Code} tab, Ctrl+Enter; in \term{Waveforms} keep \code{t=1} and make sure
the period of \code{q} is twice the period of \code{clk} (as in the figure
above); then \code{t=0}---the output freezes.

\begin{diybox}[T and JK in Verilog: Complete the Module]
\textbf{Task.} \stars{2} Describe a JK flip-flop: $J=1,K=0$ is set,
$J=0,K=1$ is reset, $J=K=1$ is toggle, $J=K=0$ is hold. \stars{3} Obtain a T
flip-flop from your JK flip-flop by connecting $J=K=T$.\\
\textbf{Expected result:} at $J=K=1$ the output changes at every edge.
\end{diybox}

\textbf{Sample answer.}
\begin{vcode}
module jk_ff (input clk, j, k, output reg q);
    always @(posedge clk) begin
        case ({j, k})                 // concatenation of two bits into one number
            2'b00: q <= q;            // hold
            2'b01: q <= 1'b0;         // reset
            2'b10: q <= 1'b1;         // set
            2'b11: q <= ~q;           // toggle
        endcase
    end
endmodule
\end{vcode}

The \term{JK flip-flop} combines the abilities of the RS and T flip-flops: $J$ is
set, $K$ is reset, $J = K = 1$ is toggle (there is no forbidden combination). In
old circuits built from the 155-series chips, JK flip-flops were very
popular. In modern programmable-logic chips only D flip-flops remain, and the
other kinds are obtained from them when needed by adding logic at the input.

\begin{table}[H]
\centering
\caption{Summary: kinds of flip-flops}
\small
\begin{tabularx}{\textwidth}{l l X}
\toprule
\textbf{Type} & \textbf{When it acts} & \textbf{What it does} \\
\midrule
RS latch & immediately & $S$ sets to 1, $R$ resets to 0, $S=R=1$ is forbidden \\
D latch  & while $C = 1$ & follows $D$, holds at $C = 0$ \\
\textbf{D flip-flop} & \textbf{on the edge} of CLK & stores $D$ at the moment of the edge \\
T flip-flop  & on the edge of CLK & at $T=1$ changes state to the opposite one \\
JK flip-flop & on the edge of CLK & $J$ is set, $K$ is reset, $J=K=1$ is toggle \\
\bottomrule
\end{tabularx}
\end{table}

\begin{ideabox}
A flip-flop stores one bit. The main memory element in digital technology is the
\textbf{D flip-flop}: it stores the value of the input $D$ at the moment of the
\emph{rising edge} of the clock signal and holds it until the next edge. For the
flip-flop to work reliably, the input must not change in the window
$t_\text{su}$--$t_\text{h}$ around the edge.
\end{ideabox}

\begin{diybox}[T and JK Flip-Flops; Project Stage]
\textbf{Task.} \stars{1} The lessons ``T flip-flop: a toggle'' and ``JK flip-flop''.
\stars{2} The lesson ``A divide-by-2 frequency divider'' (chapter ``Counters and
dividers''): make sure the output is half as fast as the clock. \stars{3}
\textbf{Project stage ``Security Panel'': remembering the alarm.} Open the
project ``An alarm with memory'' (chapter ``Projects'' of the ``Circuits and the
board'' book): the sensor \code{sensor} gave 1 once, and the output \code{y}
stays 1 until the reset \code{reset} arrives. Take as the sensor the signal
``door open and guard on'' from stage~2, and as the reset the guard's key.\\
\textbf{In Logra:} Learn (chapter ``Projects''), Designer, Waveforms, RTL.\\
\textbf{Expected result:} the door was open for one clock cycle and then closed---the
siren keeps sounding; the key clears the alarm, and the reset is stronger than the
sensor.\\
\textbf{Acceptance criterion:} all six checks of the lesson are green; Waveforms
shows a ``step'' of the output that does not disappear after the door closes; RTL
shows one flip-flop and one OR gate in the feedback path.
\end{diybox}

\textbf{The project stage in Verilog.} The same node as a module:

\vfile{\codepath/b7_flipflops/alarm_keep.v}{code/b7\_flipflops/alarm\_keep.v}

\textbf{How to read it.} Two branches in order of priority: the key resets,
otherwise the sensor sets; if neither arrives, \code{alarm} holds its previous
value (there is no \code{else} branch). The testbench \code{tb\_alarm\_keep.v}
goes through seven steps: quiet; the door opened; closed (the alarm is
remembered); the key reset it; key and sensor together (the key is stronger);
the sensor again. \textbf{How to run it in Logra:} \term{Code} tab, paste the
module, Ctrl+Enter, then in \term{Waveforms} apply \code{sensor} for one clock
cycle: \code{alarm} must stay at one until a \code{key} pulse.

\textbf{Complete it.} Add an input \code{mute} that silences the siren but does
not clear the alarm. \emph{Sample:} add a port \code{output siren} and the line
\code{assign siren = alarm \& \textasciitilde mute;}---the flip-flop itself does
not change.

\section{Common Mistakes}

\begin{warnbox}[Common Mistakes]
\begin{enumerate}
  \item \textbf{``A latch is the same as a flip-flop.''} A latch is transparent
        the whole time while $C=1$; a flip-flop takes the value only at the
        moment of the edge. Put a latch in a feedback loop (for example, $D=\overline{Q}$)---
        you get not a divide-by-2 but a noise generator.
  \item \textbf{Releasing both signals $S=R=1$ at the same time.} The state after
        that is unpredictable: both gates ``want'' to switch at once, and
        whichever happens to be slightly faster wins.
  \item \textbf{``A two-flip-flop synchronizer solves the problem forever.''} It
        only makes metastability at the output extremely
        rare, but not impossible. And it is good for one bit: a bus of
        several bits cannot be passed this way, because the bits may arrive in
        different clock cycles.
  \item \textbf{Not setting the initial state.} After power-up
        a flip-flop is in a random state. If the alarm is ``remembered'' out of
        nothing, the siren will wail right after power-on. You need a reset.
  \item \textbf{Expecting $Q$ to change at the moment of the edge.} The output changes
        $t_\text{co}$ \emph{after} the edge, and the inputs of its own
        flip-flop ``see'' the old value. That is exactly why chains of flip-flops shift
        rather than fly straight through.
\end{enumerate}
\end{warnbox}

\subsection{Common Mistakes in Verilog}

While flip-flops were circuits, mistakes were visible in the picture. When they
become code, new ones appear, and all of them come from one feature:
\emph{in the text, a combinational block and a block with memory look alike}.

\begin{warnbox}[Mistakes in Describing Memory]
\begin{enumerate}
  \item \textbf{An accidental latch from an incomplete \code{if}.} If in
        \code{always @*} the output is not assigned in all branches, the
        synthesizer inserts a latch, and Logra warns: \emph{Latch inferred}.
  \item \textbf{No \code{default}.} In a \code{case} without a \code{default}
        branch (and without a default assignment at the start of the block), the
        output holds its previous value for undescribed combinations: also a
        latch.
  \item \textbf{Blocking \code{=} in a clocked block and non-blocking \code{<=} in
        a combinational one.} In \code{always @(posedge clk)} write \code{<=}; in
        \code{always @*} write \code{=}. Breaking the rule gives a shift that
        ``flies through'' in one clock cycle (see
        Section~\ref{sec:regs-shift-v} in Chapter~\ref{ch:regs}).
  \item \textbf{No initial value after power-up.} A \code{reg} is unknown
        (\code{X}) until its first write; counters and state machines need a reset
        or an initial value in the declaration (\code{reg q = 0}).
  \item \textbf{An inconsistent reset.} \code{posedge clk or posedge rst} in the
        sensitivity list and \code{if (!rst)} inside the block: the list and the
        check must correspond to each other.
\end{enumerate}
\end{warnbox}

The first mistake is easy to see with your own eyes:

\begin{vcode}
module oops (input sel, a, output reg y);
    always @* begin
        if (sel) y = a;      // no else: y holds its value, we got a latch!
    end
endmodule
\end{vcode}

The fix is to add a branch or to assign a default value before the \code{if}:

\begin{vcode}
module oops_fixed (input sel, a, output reg y);
    always @* begin
        y = 1'b0;            // default value
        if (sel) y = a;
    end
endmodule
\end{vcode}

\section{Questions and Problems}
\label{sec:tasks-ff}

Answers are collected in Appendix~\ref{app:answers}.

\begin{taskbox}[Chapter Problems]
\begin{enumerate}
  \item \stars{1}~For an RS flip-flop made of NOR gates, which input combination is forbidden and why?
  \item \stars{2}~How does a D latch differ from a D flip-flop? Draw the response of both to the same input signal during a clock pulse.
  \item \stars{2}~The setup time is 0.4~ns, the flip-flop output delay is 0.5~ns, and the combinational logic between flip-flops is 3~ns. What is the minimum period and what is the maximum frequency?
  \item \stars{1}~How do you make a T flip-flop from a D flip-flop? What appears at the output if $T=1$ constantly?
  \item \stars{1}~What does a JK flip-flop do at $J=K=1$ on a clock edge?
  \item \stars{2}~What is metastability and in what situation does it arise? Why is an asynchronous signal passed through two flip-flops in a row?
  \item \stars{2}~\textbf{Find the error.} A student wanted to make a divide-by-2 frequency divider and connected the output $\overline{Q}$ of a D latch to its input $D$, using a 1~MHz signal as the clock. What will happen and how can it be fixed?
  \item \stars{2}~\textbf{Find the error.} In the ``alarm with memory'' the input $D$ of the flip-flop was connected only to the sensor (without feedback through an OR). What will happen after the door closes?
  \item \stars{3}~\textbf{Project.} In the Designer, build the ``alarm with memory'' without the ready-made lesson: sensor \code{sensor}, key \code{key} (reset), a synchronous flip-flop. Acceptance criterion: in Waveforms the door is open for one clock cycle, the alarm holds, the key resets, and when \code{key} and \code{sensor} are active together the key wins.
  \item \stars{3}~\textbf{Project.} One pulse per press: put a synchronizer and an edge detector in front of the alarm. Acceptance criterion: the button is held for 10 cycles, and the event counter (one press) increased by 1.
\end{enumerate}
\end{taskbox}

\section{What's Next}

We have learned to store \emph{one} bit and know when it can be trusted. But the
lock code, the alarm time and the zone number are numbers, not single bits. The
next question: what happens if we put several flip-flops side by side, and how
can we make them count time for the siren's ``beep''? That is the subject of the
next chapter: registers, shifting and counters.
